\section{Introduction}

The observable in this paper is the field strength. We regard it as a real
tempered two-form on $\R^3$, and write $F(f)$ for its pairing with a test
two-form $f$. The regularization uses a noncompact gauge field and
four-component chamber fermions, with all flux sectors included.
The order of limits is part of the construction: the volume tends to
infinity at a fixed lattice spacing, and the spacing then tends to zero.

For a lattice curvature $u$, the physical field $F^\sharp$ is the
centered plaquette step interpolation of $\eta^{-3/2}u$ at spacing
$\eta$; see \cref{sec:regularization,sec:continuum-limits} for the
finite law and interpolation. A volume-first limit is a weak limit
of fixed-spacing infinite-volume laws as $\eta\to0$. We take these
weak limits on the Hilbert space $H_{-4}$ of tempered two-forms
defined in \cref{sec:probability}. In the following statement
$\gamma_j=\operatorname{diag}(\sigma_j,-\sigma_j)$, where $\sigma_j$
are the Pauli matrices, and $d^*$ is the Euclidean adjoint of $d$.

\begin{theorem}[Field-strength OS theorem]\label{thm:main}
Fix $0<\rho<1$ and $0<\bar m\le1$, the chamber Wilson parameter and
reference mass parameter in physical units. Use the standard chamber weights
$(6,6\sqrt2,6\sqrt2,6)$. There are an odd scale factor $L\ge5$,
numbers $e_{\mathrm{phys}},\sigma>0$, and real volume-independent bare mass shifts
$m_{0,N}(e)$ for $N\in\Z_{\ge0}$ and $0<e\le e_{\mathrm{phys}}$, with
$|m_{0,N}(e)|<1/8$, with the following property.

Fix a real $\beta\in C_c^\infty(\R^3;\Lambda^2)$ such that
$f=dd^*\beta$ has
\begin{equation}\label{eq:positive-coefficient}
 b(f):=-6\tr\bigl[(S_{\bar m}\,i\gamma\cdot d^*\beta)^4\bigr]>0,
 \qquad S_{\bar m}=(\bar m+\gamma\cdot\partial)^{-1}.
\end{equation}
Such probes exist. There is $e_\beta\in(0,e_{\mathrm{phys}}]$ such that, for $0<e\le e_\beta$,
the all-flux chamber laws on the coordinate cubic tori of unit
side $n=R/\eta$, with periodic gauge fields and fermions
antiperiodic in all three coordinate directions, and with
\begin{equation}\label{eq:cutoff-scales}
 \eta=L^{-N},\quad R=8hL^T,\quad
 h=L^{\lfloor\log_L\bar m^{-1}\rfloor},\quad
 q=e\eta^{1/2},\quad \boldsymbol\mu_N=\bar m\eta+m_{0,N}(e)
\end{equation}
with $N,T\in\Z_{\ge0}$ have volume-first limit points.
Every such limit point is carried by closed tempered two-forms,
satisfies Gaussian domination with constant one, and has tensor
Schwinger functions satisfying \textup{(E0$'$)}--\textup{(E3)} below.
Let $\mu$ be any such limit point and
$\mu=\int\mu_\xi\,\dd\nu(\xi)$ its decomposition into
translation-ergodic laws. For $\nu$-almost every $\xi$, $\mu_\xi$ is
carried by closed tempered two-forms and satisfies
\begin{equation}\label{eq:main-domination}
 \int e^{F(g)}\,\dd\mu_\xi(F)\le e^{\|g\|_{L^2}^2/2}
 \qquad(g\in\Sspace(\R^3;\Lambda^2)\text{ real}).
\end{equation}
Its Schwinger functions
$S_n^\xi(g_1\otimes\cdots\otimes g_n)
=\int\prod_{a=1}^nF(g_a)\,\dd\mu_\xi(F)$
extend to the complex tensor Schwartz spaces
$\Sspace_n=\Sspace(\R^{3n};(\Lambda^2\C^3)^{\otimes n})$
and obey the following conditions, with $S_0^\xi=1$.
\begin{enumerate}
\item[\textup{(E0$'$)}] For $n\ge1$ and $\Psi\in\Sspace_n$,
\begin{equation}\label{eq:main-growth}
 |S_n^\xi(\Psi)|\le2\,66^n\sqrt{n!}\,q_{6n}(\Psi),\qquad
 q_m(\Psi)=\max_{|\alpha|\le m}\sup_y
 (1+|y|^2)^{m/2}|\partial^\alpha\Psi(y)|,
\end{equation}
where the maximum is also over tensor components.
\item[\textup{(E1)}] For every Euclidean motion $g(x)=Qx+a$,
$Q\in O(3)$, $S_n^\xi(g\cdot\Psi)=S_n^\xi(\Psi)$, where
\[
 (g\cdot\Psi)(x_1,\ldots,x_n)
 =(\Lambda^2Q)^{\otimes n}\Psi(g^{-1}x_1,\ldots,g^{-1}x_n).
\]
\item[\textup{(E2)}] Let $r$ be reflection in a plane through the origin
and $H$ either closed half-space bounded by that plane.
If $\Psi_0\in\C$ and $\Psi_k\in\Sspace_k$, $1\le k\le M$,
have $\supp\Psi_k\subset H^k$, then
\begin{equation}\label{eq:main-rp}
 \sum_{k,l=0}^M S_{k+l}^\xi
 \bigl(\Theta_k\Psi_k\otimes\Psi_l\bigr)\ge0,
 \qquad \Theta_k\Psi_k=\overline{r\cdot\Psi_k},\quad
 \Theta_0c=\overline c.
\end{equation}
\item[\textup{(E3)}] $S_n^\xi$ is unchanged by simultaneous permutation
of the arguments and their tensor indices.
\item[\textup{(E4)}] For $\Psi\in\Sspace_k$ and $\Phi\in\Sspace_l$,
\begin{equation}\label{eq:main-cluster}
 S_{k+l}^\xi(\Psi\otimes\tau_a^{\otimes l}\Phi)
 \longrightarrow S_k^\xi(\Psi)S_l^\xi(\Phi)
 \quad\text{as }|a|\longrightarrow\infty,
\end{equation}
where
\[
 (\tau_a^{\otimes l}\Phi)(x_1,\ldots,x_l)
 =\Phi(x_1-a,\ldots,x_l-a).
\]
\end{enumerate}
Moreover, for a constant $C_f$ depending only on the fixed
construction, reference mass and probe, and independent of
$0<e\le e_\beta$, the limit point and $\xi$,
\begin{equation}\label{eq:main-cumulant}
 \bigl|\kappa_4^{\mu_\xi}(f)-e^4b(f)\bigr|
 \le C_f e^{4+\sigma},\qquad
 \kappa_4^{\mu_\xi}(f)\ge\tfrac12e^4b(f)>0.
\end{equation}
In particular, $\mu_\xi$ is non-Gaussian.
The full-measure component set is asserted for each fixed parent
law, charge and probe.
\end{theorem}

Here $\kappa_4^\lambda(f)$ is the fourth cumulant of $F(f)$ under a
law $\lambda$. For the Gaussian-dominated laws considered here,
$F(f)$ has mean zero, so
$\kappa_4^\lambda(f)=S_4^\lambda(f^{\otimes4})
-3S_2^\lambda(f^{\otimes2})^2$.
The charge threshold may depend on the fixed probe.
The factorization in \textup{(E4)} is a conclusion for the
translation-ergodic components.

The tensor action in \textup{(E1)} is essential to
\textup{(E2)}. For reflection in $x_1=0$, the components $F_{12}$ and
$F_{13}$ change sign, while $F_{23}$ does not. Replacing this action
by scalar reflection changes the positivity assertion.
The growth condition is stated on all tensor tests, including tests
meeting the diagonals; bounds on products of separately smeared fields
alone do not state \textup{(E0$'$)}.


\subsection{The proof and its relation to earlier work}

\Cref{thm:main} states the tensor form of the axioms of
Osterwalder and Schrader \cite{OS1973,OS1975}. Uniform Gaussian
estimates give tightness of the cutoff laws and convergence of their
moments along weakly convergent sequences. The lattice reflection
argument supplies positivity at the
three coordinate and six diagonal mirror directions; the comparison
with mirror-adapted tori requires taking the volume limit first.
These properties pass to every prescribed volume-first limit by
\cite{ChamberLattice,ContinuumLimits}. Reflection-positive lattice
gauge theories and their transfer matrices were developed by
Osterwalder and Seiler \cite{OsterwalderSeiler1978}; the chamber
crossings and the comparison of the two torus geometries are the
construction-specific inputs used here.

The nine mirror directions give the spectral inequalities
$H_j\ge|\mathsf P_i|$ for $i\ne j$. Together with Gaussian
domination, they bound the angular frequencies of the Schwinger
functions. Closedness gives the required tensor weight at two
points; reflection positivity and the regular-polygon argument then
give that weight at every order. The proof in \cref{sec:rotations}
also treats coincidences and excludes contact terms. We apply it
to almost every ergodic component, after obtaining the Gaussian
bound and the nine reflection-positive forms on a common component
set from \cite{ErgodicComponents}. The same spectral inequalities,
together with the local correlation estimate, give clustering in
each component. The supported-tensor argument
in \cref{sec:tensor-axioms} extends these conclusions to the tests
in \textup{(E0$'$)}--\textup{(E4)}, including closed-half-space
supports in reflection positivity.

The renormalization estimates select the volume-independent mass
and control the fourth-order response. Dimock \cite{DimockUVS2022}
proves ultraviolet stability at weak coupling in a fixed finite
volume using Ba{\l}aban's renormalization group. The block analysis
used here adapts this framework to chamber fermions and two
reference tori. Decay estimates for the local inverses control the
polymer activities; extraction of the vacuum and mass terms leaves
a normalized interaction that contracts under block averaging.
The nonlinear step and the common tuning theorem give controlled
trajectories for one mass and vacuum assignment on both reference
tori \cite{NormalizedContraction,NonlinearStep,CoverTuning}.
Terminal pressure estimates then enter the quantitative response
bound \cite{TerminalPressure,ReferenceResponse}. This control of
the interacting flow supplies the non-Gaussianity estimate, while
the Gaussian bounds supply the compactness and moment bounds used
in the limit passage.

The transfer from the reference tori uses chessboard estimates,
whose general framework is due to Fr\"ohlich, Israel, Lieb, and
Simon \cite{FILS1978}. The volume-transfer theorem of
\cite{VolumeTransfer} controls products of translated Wick
powers as well as the fourth cumulant. For a prescribed parent law,
these product bounds control invariant conditional expectations
and yield the cumulant estimate in almost every ergodic component
of that same law. The free fourth coefficient converges to $b(f)$,
and the continuum calculation in \cite{PositiveProbe} provides a
compact probe with $b(f)>0$. Intersecting the resulting full-measure
component set with the structural one, and reducing the
probe-dependent charge interval, gives \eqref{eq:main-cumulant}
along with all five tensor axioms.
The proof of \cref{thm:main} is completed in
\cref{sec:main-theorem-closure}.
